\section{Toward Integrated Open-World Agents}
\label{sec:perspectives_integration}

The preceding sections reviewed the state of the field across perception, planning, learning, execution, and the emerging role of foundation models. A recurring theme is fragmentation: methods that address different aspects of the open-world problem are typically developed and evaluated in isolation from one another. This section briefly discusses how the contributions of this dissertation relate to the broader goal of integrated open-world agents.

The three contributions of this dissertation---benchmark environments (\cref{ch:benchmarks}), action abstractions for adaptation (\cref{ch:rapid_learn}), and interaction abstractions for generalization (\cref{ch:generalization})---address complementary aspects of the open-world problem. The benchmarks provide controlled settings for evaluating agents under diverse classes of novelty, measuring both the dynamics of performance degradation and the efficiency of recovery. The action abstractions provide the mechanism for efficient adaptation when the agent's symbolic model becomes incorrect, converting global failures into targeted learning problems. The interaction abstractions provide representations that support generalization across objects and embodiments, reducing the scope of what must be relearned when conditions change.

These contributions operate at different levels of the agent architecture, and their integration represents a natural next step. As discussed in the summary of \cref{ch:generalization}, the force-space policies developed as interaction abstractions could serve as executors within the hybrid planning--learning framework that leverages action abstractions. In such an integrated system, the planning layer would provide task decomposition and failure localization; the interaction abstractions would provide generalizable execution; and when novelty renders an executor ineffective, the action abstraction framework would focus learning on acquiring a replacement. The benchmarks would provide the evaluation infrastructure for measuring whether such an integrated system achieves the adaptation and generalization capabilities that neither component provides alone.

Realizing this integration fully requires addressing the perception gap identified in \cref{sec:rw_sensing}: the current contributions assume access to task-relevant state information, and connecting open-world perception to the downstream reasoning and execution pipeline remains an important open challenge. It also requires scaling the symbolic planning components to richer domains, and developing the formal connections between the novelty taxonomy of \cref{sec:novelty_taxonomy} and the adaptation mechanisms to enable more principled selection of recovery strategies based on novelty type. These directions are discussed further in \cref{ch:conclusion}.